\section{Una montaña infinita: los problemas son inevitables, pero tienen solución}

Esta sección parte de «The Beginning of Infinity». Yang Zhilin (杨植麟) dice que los problemas son inevitables, pero que tienen solución; quizá esta montaña nevada no tenga cima, pero él espera que nunca la tenga. La metáfora describe bien a una empresa de modelos: cada vez que mejora la capacidad del modelo se desbloquean nuevos escenarios y, al mismo tiempo, surgen nuevos problemas. El preentrenamiento, el RLHF, el pensamiento largo, las capacidades Agentic, el uso de herramientas y los sistemas de evaluación son como rutas a distintas alturas de la montaña nevada.

\begin{figure}[H]
\centering
\includegraphics[width=0.92\textwidth]{infinite-mountain-map.png}
\caption{Una montaña infinita: los problemas son inevitables, pero tienen solución; el avance de los modelos abre continuamente nuevos problemas. Diagrama conceptual propio, elaborado a partir de la conversación de 00:00:02--00:05:18.}
\end{figure}

\begin{knowledgebox}{Lectura de la figura: lo infinito no es pesimismo, sino el oxígeno de la investigación}
Si la montaña tuviera cima, la empresa de modelos pronto se convertiría en una empresa puramente de ingeniería; si no la tiene, la investigación, el producto y la organización seguirán enfrentando nuevos problemas. La montaña infinita significa que el espacio de problemas de largo plazo sigue existiendo.
\end{knowledgebox}

\subsection{Test-time scaling: el pensamiento largo y los Agent apuntan en la misma dirección}

La sección anterior trató de que «los problemas seguirán surgiendo»; esta entra en el primer juicio técnico. Yang Zhilin considera que tanto el aprendizaje por refuerzo basado en pensamiento largo como el aprendizaje por refuerzo de Agent apuntan hacia el test-time scaling. El test-time scaling consiste en invertir más cómputo en la etapa de inferencia/prueba: cadenas de razonamiento más largas, más llamadas a herramientas, más verificación y más iteraciones.

\begin{figure}[H]
\centering
\includegraphics[width=0.96\textwidth]{test-time-scaling-map.png}
\caption{Test-time Scaling: tanto el RL de pensamiento largo como el RL de Agent amplían el cómputo en tiempo de inferencia. Diagrama conceptual propio, elaborado a partir de la conversación de 00:03:38--00:24:58.}
\end{figure}

\begin{knowledgebox}{Lectura de la figura: el escalamiento en tiempo de prueba no consiste solo en generar más token}
El pensamiento largo hace que el modelo gaste más token en el razonamiento interno; los Agent permiten que el modelo obtenga retroalimentación externa a través de herramientas y del entorno. Ambos gastan más cómputo en la etapa de prueba; uno se inclina por el pensamiento interno y el otro por la acción externa.
\end{knowledgebox}

\begin{importantbox}{Fórmula intuitiva del test-time scaling}
\[
\text{Capability} \approx f(\text{model}, \text{test-time compute}, \text{feedback})
\]
Donde \(\text{model}\) es la capacidad del modelo base, \(\text{test-time compute}\) es la inversión adicional en la etapa de inferencia y \(\text{feedback}\) es la verificación, lo que devuelven las herramientas o la retroalimentación del entorno. Lo decisivo en un Agentic LLM es convertir estos dos últimos elementos en un sistema.
\end{importantbox}

\subsection{De L1 a L5 no necesariamente en serie}

La sección anterior descompuso el test-time scaling en pensamiento interno y acción externa; esta explica además que estas capacidades no necesariamente suben los peldaños de forma lineal. En la entrevista se menciona que de L1 a L5 no hay necesariamente una relación en serie: Claude logró muy buenos resultados en Agent aun cuando su reasoning no era necesariamente el más fuerte. Este juicio es importante: el desarrollo de las capacidades de los modelos no es una escalera fija; se puede apostar en paralelo por Reasoning, Agent, herramientas, producto y retroalimentación. Solo cuando el modelo participa en el proceso de desarrollo es posible desbloquear la verdadera etapa de Innovator.

\begin{figure}[H]
\centering
\includegraphics[width=0.96\textwidth]{l1-l5-agent-ladder.png}
\caption{De L1 a L5 no necesariamente en serie: se puede apostar en paralelo por Reasoning y Agent; Claude es un ejemplo típico. Diagrama conceptual propio, elaborado a partir de la conversación de 00:12:00--00:24:58.}
\end{figure}

\begin{warningbox}{No hay que tratar los niveles de capacidad como una escalera mecánica}
Si se entiende que de L1 a L5 hay que superar cada nivel en orden, se subestiman las rutas de producto y de sistema. La capacidad de Agent puede despuntar primero en ciertas tareas y, después, impulsar a su vez el entrenamiento del modelo y el ecosistema de herramientas.
\end{warningbox}

\begin{knowledgebox}{Explicación didáctica de las etapas L}
\begin{tabularx}{\textwidth}{p{0.16\textwidth}p{0.34\textwidth}X}
\toprule
Etapa & Intuición & Variables decisivas \\
\midrule
L1/L2 & Chat, preguntas y respuestas, reasoning básico & Base del modelo y alineación. \\
L3 & Puede usar herramientas para completar tareas bien definidas & Interfaces de herramientas, definición de tareas, retroalimentación. \\
L4 & Participa en el desarrollo, la innovación y la resolución de problemas complejos & Agent de largo alcance, verificación, capacidad de programación. \\
L5 & Sistemas con mayor autonomía & Seguridad, organización, permisos y retroalimentación continua. \\
\bottomrule
\end{tabularx}
\end{knowledgebox}

\subsection{Resumen de la sección}

«La montaña infinita» da a este episodio su tono de fondo: cada paso que avanza la capacidad del modelo abre nuevos problemas; el test-time scaling, por su parte, marca la dirección técnica: hacer que el modelo piense más, actúe más y obtenga más retroalimentación durante la inferencia.
